% =============================================================================
\section{px4\_link.hpp: lo que comparten todas}
% =============================================================================

\begin{frame}{Por qué una clase común}
  \begin{columns}[c]
    \column{0.5\textwidth}\small
    En \cmd{px4\_examples} cada nodo repetía \textasciitilde80 líneas:
    suscripciones, publishers, \cmd{send\_command()}, \cmd{now\_us()}\ldots
    \vspace{0.4em}

    Aquí eso vive en \cmd{x500::Px4Link} y cada misión se queda solo con
    \textbf{su máquina de estados}.
    \vspace{0.4em}
    \begin{idea}
      \small Si PX4 cambia la versión de un mensaje (\cmd{vehicle\_status\_v1}\ldots),
      se corrige en un solo archivo.
    \end{idea}
    \column{0.47\textwidth}
    \begin{tikzpicture}[every node/.style={font=\scriptsize}]
      \node[draw=azul,fill=azulclaro,rounded corners,minimum width=2.4cm,minimum height=1.6cm] (l) {Px4Link};
      \node[draw,rounded corners,left=1.2cm of l,minimum width=1.4cm] (p) {PX4};
      \node[draw=azul,rounded corners,above=0.7cm of l,minimum width=2.4cm] (n) {tu misión};
      \draw[-{Stealth}] ([yshift=0.4cm]p.east) -- node[above,font=\tiny] {/fmu/out} ([yshift=0.4cm]l.west);
      \draw[{Stealth}-] ([yshift=-0.4cm]p.east) -- node[below,font=\tiny] {/fmu/in} ([yshift=-0.4cm]l.west);
      \draw[{Stealth}-{Stealth}] (n) -- (l);
    \end{tikzpicture}
  \end{columns}
\end{frame}

\begin{frame}[fragile]{Uso}
  \begin{lstlisting}[style=cpp,caption={Dentro de cualquier rclcpp::Node}]
#include "x500_retos/px4_link.hpp"

class MiMision : public rclcpp::Node {
  x500::Px4Link px4_{this};         // crea suscripciones y publishers

  void loop() {
    const auto & l = px4_.local();  // posición NED (m), velocidad, heading
    if (px4_.armed() && px4_.in_offboard()) { ... }
    px4_.offboard(setpoint);        // heartbeat + setpoint, CADA ciclo
    px4_.command(VehicleCommand::VEHICLE_CMD_NAV_LAND);
  }
};
  \end{lstlisting}
\end{frame}

\begin{frame}{Lo que lee de PX4}
  \centering\small
  \begin{tabular}{lll}
    \toprule
    Método & Topic & Para qué \\
    \midrule
    \cmd{status()} & \cmd{/fmu/out/vehicle\_status\_v1} & armado, modo (nav\_state) \\
    \cmd{local()} & \cmd{/fmu/out/vehicle\_local\_position\_v1} & x, y, z, vx, vy, vz, heading \\
    \cmd{global()} & \cmd{/fmu/out/vehicle\_global\_position} & lat, lon, alt (AMSL) \\
    \cmd{home()} & \cmd{/fmu/out/home\_position\_v1} & dónde está ``casa'' \\
    \cmd{attitude()} & \cmd{/fmu/out/vehicle\_attitude} & cuaternión q = \{w, x, y, z\} \\
    \bottomrule
  \end{tabular}
  \vspace{0.6em}

  \raggedright Y atajos: \cmd{armed()}, \cmd{in\_offboard()}, \cmd{in\_hold()}, \cmd{local\_ready()},
  \cmd{gps\_ready()}, \cmd{distance\_to(p)}.
  \vspace{0.3em}
  \begin{idea}
    \small También escucha \cmd{vehicle\_command\_ack} y avisa con un \cmd{WARN} cuando PX4
    \textbf{rechaza} un comando (una vez por cada respuesta distinta).
  \end{idea}
\end{frame}

\begin{frame}[fragile]{El setpoint de offboard}
  \begin{lstlisting}[style=cpp,caption={x500::Setpoint: lo que se deja en NaN, PX4 no lo usa}]
struct Setpoint {
  std::array<float, 3> position{kNaN, kNaN, kNaN};      // m, NED
  std::array<float, 3> velocity{kNaN, kNaN, kNaN};      // m/s
  std::array<float, 3> acceleration{kNaN, kNaN, kNaN};  // m/s^2
  float yaw = kNaN;        // rad, 0 = norte, + hacia el este
  float yawspeed = kNaN;   // rad/s, feedforward de giro
};
  \end{lstlisting}
  \small \cmd{offboard(sp)} publica dos cosas, cada vez:
  \begin{enumerate}
    \item \cmd{OffboardControlMode}: qué campos traen valor (el \emph{heartbeat}).
    \item \cmd{TrajectorySetpoint}: los valores.
  \end{enumerate}
\end{frame}

\begin{frame}{Setpoints mixtos}
  \small Basta un eje con valor para que ese tipo cuente. Así se combinan:
  \vspace{0.3em}
  \centering
  \begin{tabular}{lll}
    \toprule
    position & velocity & Qué hace PX4 \\
    \midrule
    $\{N, E, Z\}$ & NaN & ir a ese punto \\
    $\{N, E, Z\}$ & $\{v_N, v_E, v_Z\}$ & ir al punto usando la velocidad como \emph{feedforward} \\
    $\{\text{NaN}, \text{NaN}, Z\}$ & $\{v_N, v_E, \text{NaN}\}$ & avanzar por velocidad a altura fija (cueva, vuelo\_prueba) \\
    $\{N, E, \text{NaN}\}$ & $\{\text{NaN}, \text{NaN}, v_Z\}$ & quedarse en $(N, E)$ bajando a $v_Z$ (aterrizaje de precisión) \\
    \bottomrule
  \end{tabular}
\end{frame}

\begin{frame}{Feedforward: no ir ``atrás'' del setpoint}
  \begin{columns}[c]
    \column{0.5\textwidth}
    \begin{tikzpicture}[scale=0.9]
      \draw[gray!50] (0,0) circle (1.6);
      \draw[-{Stealth},azul,thick] (1.6,0) arc (0:70:1.6);
      \fill[azul] (70:1.6) circle (2pt) node[above right,font=\tiny] {setpoint};
      \fill[rojo] (35:1.6) circle (2pt) node[right,font=\tiny] {dron (solo posición)};
      \fill[verde] (64:1.6) circle (2pt) node[left,font=\tiny,xshift=-2pt] {con feedforward};
    \end{tikzpicture}
    \column{0.48\textwidth}\small
    \begin{itemize}
      \item Con \textbf{solo posición}, PX4 corrige \emph{después} de ver el error: en una
        trayectoria que se mueve, siempre va atrás.
      \item Con la \textbf{velocidad} (y aceleración) de la trayectoria, ya sabe cómo moverse:
        el error queda casi en cero.
      \item Lo mismo con el yaw: \cmd{yawspeed} le dice qué tan rápido girar.
    \end{itemize}
  \end{columns}
\end{frame}

\begin{frame}{El heartbeat de offboard}
  \centering
  \begin{tikzpicture}[x=1.1cm]
    \draw[-{Stealth}] (0,0) -- (10.5,0) node[right,font=\scriptsize] {t};
    \foreach \x in {0.25,0.5,...,10} \draw[azul] (\x,0) -- (\x,0.3);
    \node[font=\tiny,azul,above] at (1.25,0.3) {setpoints a 20 Hz};
    \draw[-{Stealth},thick,naranja] (1,-0.9) -- (1,-0.05);
    \node[font=\tiny,below,align=center] at (1,-0.9) {$\approx$1 s: pedir\\OFFBOARD + armar};
    \draw[-{Stealth},thick,verde] (2.2,-0.9) -- (2.2,-0.05);
    \node[font=\tiny,below,align=center] at (2.6,-0.9) {PX4 entra\\a offboard};
    \draw[rojo,thick] (7.5,-0.1) -- (7.5,0.6);
    \node[font=\tiny,rojo,align=center] at (8.6,0.9) {el nodo se cae: \textgreater0.5 s sin\\heartbeat = failsafe};
  \end{tikzpicture}
  \vspace{0.6em}
  \begin{itemize}\small
    \item PX4 \textbf{no entra} a offboard si no está recibiendo setpoints ya.
    \item Mientras el nodo controla, hay que publicar \textbf{en cada ciclo} (todos usan un timer de 50 ms).
    \item Si se deja de publicar, PX4 entra en failsafe: eso también se prueba a propósito.
  \end{itemize}
\end{frame}

\begin{frame}{Comandos que usamos}
  \centering\scriptsize
  \begin{tabular}{lll}
    \toprule
    Comando & Parámetros & Atajo en Px4Link \\
    \midrule
    \cmd{DO\_SET\_MODE} & p1 = 1, p2 = 6 (OFFBOARD) & \cmd{request\_offboard\_and\_arm()} \\
    \cmd{DO\_SET\_MODE} & p1 = 1, p2 = 4, p3 = 3 (Hold) & \cmd{request\_hold()} \\
    \cmd{COMPONENT\_ARM\_DISARM} & p1 = 1 & \cmd{arm()} \\
    \cmd{NAV\_LAND} & --- & \cmd{land()} \\
    \cmd{NAV\_RETURN\_TO\_LAUNCH} & --- & \cmd{return\_home()} \\
    \cmd{NAV\_TAKEOFF} & p5, p6 = lat, lon; p7 = alt AMSL & \cmd{command(...)} \\
    \cmd{DO\_REPOSITION} & p1 = vel, p2 = 1, p5..p7 = lat, lon, alt & \cmd{command(...)} \\
    \cmd{DO\_SET\_HOME} & p1 = 0, p5..p7 = lat, lon, alt & \cmd{command(...)} \\
    \bottomrule
  \end{tabular}
  \vspace{0.5em}

  \raggedright\footnotesize \textbf{Hold antes de armar}: desde Land (como queda al aterrizar) PX4 no deja armar.
  Por eso \cmd{entrega\_domicilio} pasa por Hold antes de volver a despegar.
\end{frame}

\begin{frame}{El reloj de PX4, no el de la pared}
  \begin{columns}[c]
    \column{0.55\textwidth}\small
    PX4 SITL corre en \emph{lockstep} con Gazebo: si la simulación va lenta, PX4 también.
    Pero \cmd{now()} de ROS es el reloj de la pared.
    \vspace{0.3em}

    En una prueba el mundo del puente corría a \textbf{0.27$\times$}: la hélice ``avanzaba''
    casi 4 veces más rápido de lo que el dron podía seguir y el yaw se quedaba 40° atrás.
    \vspace{0.3em}
    \begin{idea}
      \small \cmd{px4\_.clock()}: el timestamp de la posición de PX4, en segundos. Úsenlo para
      todo lo que dependa del tiempo de la misión.
    \end{idea}
    \column{0.42\textwidth}
    \begin{tikzpicture}[x=0.55cm,y=0.5cm]
      \draw[-{Stealth}] (0,0) -- (7,0) node[right,font=\tiny] {pared (s)};
      \draw[-{Stealth}] (0,0) -- (0,5) node[above,font=\tiny] {simulación (s)};
      \draw[verde,thick] (0,0) -- (6,4.5) node[right,font=\tiny] {RTF 1.0};
      \draw[rojo,thick] (0,0) -- (6,1.6) node[right,font=\tiny] {RTF 0.27};
    \end{tikzpicture}
  \end{columns}
\end{frame}
